\paragraph*{Четыре свойства характеристических функций с доказательством}

\begin{enumerate}
    \item $|\charac{\xi}{t}| \leqslant \charac{\xi}{0} = 1$
    
    \begin{proof} $|\charac{\xi}{t}| = |\E e^{it\xi}| \leqslant \E|e^{it\xi}| = \E(1) = 1 = \charac{\xi}{0}$\end{proof}
    
    \item $\charac{\xi}{t}$ равномерно непрерывна на $\R$
    
    \begin{proof} Пусть $\eps > 0$. Так как $\sin(t\xi) \to 0$ почти наверное при $t \to 0$, по теореме Лебега о мажорируемой сходимости $\E|\sin(t\xi)| \to 0$ при $t \to 0$
    $$
    \exists \delta > 0: \forall t \in (-\delta, \delta) \ \E|\sin(t\xi)| < \frac{\eps}{4}
    $$
    Пусть $|s - t| < \delta$, тогда
    
   $|\charac{\xi}{t} - \charac{\xi}{s}| = |\E e^{is\xi}\brackets{e^{i(t - s)\xi} - 1}| \leqslant \E|e^{i(t - s)\xi} - 1| =$ \\ $ \E\sqrt{(\cos((t - s)\xi) - 1)^2 + \sin^2((t - s)\xi)} =
 \E\sqrt{2 - 2\cos((t - s)\xi)} = 2\E\Biggl|\sin\brackets{\frac{t - s}{2}\xi}\Biggr| < \eps
    $ \ \end{proof}
    
    \item $\forall t \in \R \ \charac{\xi}{t} \in \R \Longleftrightarrow \xi = -\xi$ по распределению.
    
    \begin{proof} 
    $\Longrightarrow$ 
    $\forall t \in \R \ \charac{\xi}{t} \in \R$, тогда
    $$
    \charac{\xi}{t} = \E e^{it\xi} = \E \brackets{\cos(t\xi) + i\sin(t\xi)} = \E \brackets{\cos(t\xi)} = \E \brackets{\cos(t\xi) - i\sin(t\xi)} = \charac{\xi}{-t}
    $$
    
    $\Longleftarrow$
    $\xi = -\xi$ по распределению, тогда $\E\sin(t\xi) = \E\sin(-t\xi) = -\E\sin(t\xi) = 0, \varphi_\xi(t) = \E\cos(t\xi)$  \end{proof}
     
    \item $\forall t \in \R \ \charac{\xi}{-t} = \overline{\charac{\xi}{t}}$
    
    \begin{proof} $\charac{\xi}{-t} = \E\brackets{\cos(-t\xi) + i\sin(-t\xi)} = \E\brackets{\cos(t\xi) - i\sin(t\xi)} = \E \overline{e^{it\xi}}  =\overline{\charac{\xi}{t}}$  \end{proof}
\end{enumerate}

\begin{theorem}[Об обращении (б/д)]  Пусть $\charac{\xi}{x}$ - характеристическая функция $\xi$. \end{theorem} 

\begin{itemize}
    \item Для любых двух точек непрерывности $a < b$ функции распределения $F_\xi$ верно, что  
    $$
    F_\xi(b) - F_\xi(a) = \frac{1}{2\pi}\lim\limits_{c \to \infty}\int\limits_{-c}^c \frac{e^{-ita} - e^{-itb}}{it} \charac{\xi}{t}dt
    $$
    \item Если $\int\limits_\R |\charac{\xi}{t}| dt < \infty$, то $\xi$-абс. непр. с $p_\xi(x) = \frac{1}{2\pi} \int\limits_\R e^{-itx} \charac{\xi}{t} dt$
\end{itemize}

\begin{lemmanote} 

Обратим внимание, что формулы выше по сути являются преобразованием Фурье (либо прямым, либо обратным).

\end{lemmanote} 

\paragraph*{Теорема Бохнера-Хинчина}

\begin{definition}  $\varphi: \R \to \mathbb{C}$ неотрицательно определена, если 
$$
\forall n \in \N \ \ \forall t_1, \ldots, t_n \in \R \ \  \forall z_1, \ldots, z_n \in \mathbb{C} \sum\limits_{j, k = 1}^n \varphi(t_j - t_k)z_j\overline{z_k} \geqslant 0
$$ \end{definition} 

\begin{theorem}[Бохнера-Хинчина]  

Если $\varphi$ непрерывна и $\varphi(0) = 1$, то 
\[\varphi - \text{хар. функция} \Leftrightarrow \varphi \text{ неотрицательно определена.}\] 

\end{theorem} 

\begin{proof} 
$\Longleftarrow$ Без доказательства

$\Longrightarrow$ Пусть $\varphi$ - характеристическая, тогда

$   \sum\limits_{j, k = 1}^n \varphi(t_j - t_k)z_j\overline{z_k} = \sum\limits_{j, k = 1}^n \E e^{i(t_j - t_k)\xi}z_j\overline{z_k} = \sum\limits_{j, k = 1}^n \E \brackets{e^{it_j\xi} \cdot \overline{e^{it_k\xi}}}z_j\overline{z_k} =$ \\
    $= \E\brackets{\sum\limits_{j, k = 1}^n z_j e^{it_j\xi} \cdot \overline{z_k e^{it_k\xi}}} = \E\Biggl|\sum\limits_{j = 1}^n z_k e^{it_k\xi}\Biggr|^2 \geqslant 0
$
\end{proof}